\section{Cartesian targets and optical measures}\label{sec:optics}

\subsection{An unsquared optical-path construction}
Fix a vacuum wavelength $\lambda_{\rm opt}$, positive unequal homogeneous indices
$n_o,n_i$, and distinct conjugate points $\bm O,\bm I$. The signs $s_o,s_i$ are
$+1$ for a real object or real image, respectively, and $-1$ for a virtual one.
For an interface point $\bm x$ distinct from both foci, let
$R_o=|\bm x-\bm O|$ and $R_i=|\bm x-\bm I|$. Define the signed optical path
$\Psi$ and a constant path level $C$ by
\begin{equation}
 \Psi(\bm x)=s_o n_oR_o+s_i n_iR_i,\qquad
 \Gamma_{\rm opt,\star}\subset\{\bm x:\Psi(\bm x)=C\}.
 \label{eq:cartesian-implicit}
\end{equation}
Both $\Psi$ and $C$ have units of length. The constant is fixed by a chosen
vertex or by compatible mounting and volume data. For two real conjugates this
is the usual weighted sum of distances. Signed paths extend the construction to
virtual conjugates. The formulas and Cartesian-surface classification follow
Silva-Lora and Torres \cite{silva2020}; the mechanical and control formulation
below is separate from that paper's two-surface singlet construction.

Let $\bm d_o=s_o(\bm x-\bm O)/R_o$ be the incident propagation direction and
$\bm d_i=s_i(\bm I-\bm x)/R_i$ the desired transmitted direction. Then
\begin{equation}
 \nabla\Psi=n_o\bm d_o-n_i\bm d_i,\qquad
 \Gs\Psi=0\ \Longrightarrow\ n_o\Pg\bm d_o=n_i\Pg\bm d_i.
 \label{eq:fermat-snell}
\end{equation}
This directly supplies the tangential Snell condition. Squaring
\cref{eq:cartesian-implicit} gives a quartic, but can add branches; its real
roots alone do not certify the intended refracting surface.

To select transmission, let $\bm m$ be the unit surface normal pointing from the
incident medium into the transmitted medium. Thus $\bm m=\pm\nn$ depending on
which fluid light enters first. Set $c_o=\bm d_o\cdot\bm m$ and
$\eta_n=n_o/n_i$. The physically transmitted unit direction $\bm d_t$ is
\begin{equation}
 \bm d_t=\eta_n\bm d_o+
 \left[\sqrt{1-\eta_n^2(1-c_o^2)}-\eta_n c_o\right]\bm m.
 \label{eq:snell-vector}
\end{equation}
The branch must satisfy $c_o>0$, a nonnegative square-root argument,
$\bm d_t\cdot\bm m>0$, and $\bm d_t=\bm d_i$ throughout the admitted aperture.
These exclude back-facing and totally internally reflecting branches. A finite
distance from grazing transmission is also needed for well-conditioned optical
sensitivities.

\subsection{The four axial conjugate parameters}
Choose the vertex as the origin, with optical propagation there in the positive
$z$ direction. Let $z_o,z_i$ be the signed axial coordinates of the conjugates,
and define cylindrical radius $r=\sqrt{x^2+y^2}$. For finite nonzero conjugates,
$s_o=-\operatorname{sgn}(z_o)$ and $s_i=\operatorname{sgn}(z_i)$. The level fixed
by the vertex is $C=-n_oz_o+n_iz_i$. In particular, for $z_o<0<z_i$,
\begin{equation}
 n_o\sqrt{r^2+(z-z_o)^2}+n_i\sqrt{r^2+(z_i-z)^2}
 =-n_oz_o+n_iz_i.
 \label{eq:real-cartesian}
\end{equation}

For the paper's parametric representation, define inverse lengths
$a=1/z_o$, $b=1/z_i$, and $D=n_i a-n_o b$. When the denominators are nonzero,
define the form parameters $\mathsf O$ (inverse length), $\mathsf G$
(dimensionless), $\mathsf T$ (inverse length cubed), and $\mathsf S$ (inverse
length squared) as
\begin{align}
 \mathsf O&=\frac{n_i b-n_o a}{n_i-n_o}, &
 \mathsf G&=\frac{(n_i^2a-n_o^2b)^2}{n_i n_o D(n_i b-n_o a)},\label{eq:gots-a}\\
 \mathsf T&=\frac{(n_i-n_o)(n_i+n_o)^2a^2b^2}{4n_i n_oD}, &
 \mathsf S&=\frac{(n_i+n_o)(n_i^2a-n_o^2b)ab}{2n_i n_oD}.
 \label{eq:gots-b}
\end{align}
The signed vertex graph curvature is $z''(0)=\mathsf O$; it is not the total
mechanical curvature $\kappa$. For an upward physical normal at an axisymmetric
vertex, $\kappa(0)=-2\mathsf O$.

Define the polar distance $\varrho=\sqrt{r^2+z^2}$, which is \emph{not} the
cylindrical radius $r$ or either fluid density. The parameters satisfy
$\mathsf S^2=\mathsf O\mathsf G\mathsf T$. The Cartesian polynomial in these
coordinates is
\begin{equation}
 \mathsf O\mathsf Gz^2-2(1+\mathsf S\varrho^2)z
       +(\mathsf O+\mathsf T\varrho^2)\varrho^2=0.
 \label{eq:gots-polynomial}
\end{equation}
Solving this quadratic in $z$, selecting the branch through the vertex, and
rationalizing yields
\begin{equation}
 z(\varrho)=
 \frac{(\mathsf O+\mathsf T\varrho^2)\varrho^2}
 {1+\mathsf S\varrho^2+
   \sqrt{1+(2\mathsf S-\mathsf O^2\mathsf G)\varrho^2}},\qquad
 r(\varrho)=\sqrt{\varrho^2-z(\varrho)^2}.
 \label{eq:gots-parametric}
\end{equation}
The unsquared path and transmission tests remain necessary. Singular GOTS
coordinates need not imply a singular physical surface; in that event use
\cref{eq:cartesian-implicit} directly. Equal indices, coincident foci, and
degenerate conjugates require their own optical interpretation rather than
division by a vanishing coefficient.

For collimated incidence, take $z_o\to-\infty$ and subtract the divergent constant
path. The finite condition for a real image is
$n_oz+n_i(\sqrt{r^2+(z_i-z)^2}-z_i)=0$. Here $\mathsf S=\mathsf T=0$ and the
conic constant $K=-(n_o/n_i)^2$. Writing the signed vertex curvature as
$c_v=n_i/[z_i(n_i-n_o)]$ gives
\begin{equation}
 z(r)=\frac{c_v r^2}{1+\sqrt{1-(1+K)c_v^2r^2}}.
 \label{eq:collimated-conic}
\end{equation}
For finite conjugates, a Cartesian surface is generally not an exact conic, and
$\mathsf G$ must not be identified with a fitted conic constant.

\subsection{Geometric compatibility and limits of arbitrariness}
For fixed homogeneous indices and two fixed points, $\Psi$ is invariant under
rotation about their joining axis. Consequently a regular connected Cartesian
level surface belongs to a surface of revolution. An arbitrary non-axisymmetric
surface is generally incompatible with exact single-interface point-to-point
stigmatism under these optical assumptions. Arbitrary aspheres remain valid
mechanical targets, but require a correspondingly different optical objective
or incident wavefront.

More generally, let $S_o(\bm x)$ and $S_i(\bm x)$ be incident and desired
transmitted eikonal functions, measured as optical path lengths. In homogeneous
media they satisfy $|\nabla S_o|=n_o$ and $|\nabla S_i|=n_i$, with gradients
pointing along propagation. A purely refracting interface with no imposed
phase discontinuity must lie on $S_o-S_i=C$. The point-conjugate choice is
$S_o=s_on_oR_o$, $S_i=-s_in_iR_i$, recovering
\cref{eq:cartesian-implicit}. A non-axisymmetric target can thus be coupled to
a specified different wavefront-design problem, but the necessary incident
wavefront must be physically supplied rather than assumed to arise from the
same two point conjugates.

The target must also be compatible with the available fluid mass, supporting
boundaries and allowed contact-line motion. In an incompressible model, fixing
the entire interface and its other boundaries fixes the enclosed volume. A
single optical path constant cannot generally enforce an arbitrary boundary
curve and an independently prescribed volume. Possible design freedoms are a
non-optical annulus, a movable support, a reservoir, or revised conjugates; none
is silently introduced here.

The incident directions in \cref{eq:fermat-snell} are specified in a homogeneous
incident medium. If light first crosses a window or another liquid surface,
that preceding optical system must produce those directions and optical paths.
Similarly, solidifying a liquid can change its index and shape. A liquid
Cartesian diopter is therefore not automatically a Cartesian finished lens.
Exact axial stigmatism at one wavelength does not establish the Abbe sine
condition, off-axis aplanatism, or achromatism.

\subsection{Optical objectives and a slope-sensitive identity}
Let $\mathcal P$ be a two-dimensional entrance pupil, $\bm\xi\in\mathcal P$ its
ray label, and $\bm X_\Gamma(\bm\xi)$ the ray's first intended intersection with
the interface. Every admitted ray must reach the intended transmitting patch
and have an unobstructed optical continuation. Let $\mathrm d\mu(\bm\xi)$ be a nonnegative normalized measure
representing the selected illumination, so $\int_{\mathcal P}\mathrm d\mu=1$.
Define the mean signed path $\overline\Psi=\int_{\mathcal P}
\Psi(\bm X_\Gamma)\mathrm d\mu$ and path departure
$W=\Psi(\bm X_\Gamma)-\overline\Psi$. Choose a positive optical-path tolerance
$W_{\rm ref}$ with units of length and an angular tolerance
$\theta_{\rm ref}>0$ in radians. Two dimensionless objectives are
\begin{equation}
 J_W=\int_{\mathcal P}\frac{W^2}{W_{\rm ref}^2}\mathrm d\mu,
 \qquad
 J_d=\int_{\mathcal P}
 \frac{\norm{\Pg(\bm d_t-\bm d_i)}^2}{\theta_{\rm ref}^2}\mathrm d\mu.
 \label{eq:optical-objectives}
\end{equation}
The projector and directions are evaluated at $\bm X_\Gamma$. Removing
$\overline\Psi$ removes an optically irrelevant piston; it does not move either
conjugate. A specified geometric target may additionally fix $C$ and the vertex.

For any physically transmitted surface, tangential Snell refraction gives the
exact identity
\begin{equation}
 n_i\Pg(\bm d_t-\bm d_i)=\Gs\Psi.
 \label{eq:opd-slope-identity}
\end{equation}
Thus a small mean-square path error alone does not control ray error: a small
amplitude but rapidly varying path departure can have a large surface gradient.
Using $J_d$, or an equivalent derivative-sensitive optical objective, prevents
accepting such a surface solely on its path variance. Near grazing transmission,
even small tangential direction errors can produce large changes in the normal
component; the non-grazing condition must therefore be retained.

For a real image, define a detector plane through $\bm I$ with unit normal
$\bm e_i$, and require $\bm d_t\cdot\bm e_i>0$. Its ray intersection is
$\bm Y=\bm X_\Gamma+[(\bm I-\bm X_\Gamma)\cdot\bm e_i]
\bm d_t/(\bm d_t\cdot\bm e_i)$, with a positive propagation distance to that
plane. For a prescribed positive spot-radius tolerance
$r_{\rm ref}$, an additional dimensionless cost is
\begin{equation}
 J_{\rm spot}=\int_{\mathcal P}
       \frac{|\bm Y-\bm I|^2}{r_{\rm ref}^2}\mathrm d\mu.
 \label{eq:spot-cost}
\end{equation}
This keeps the requested image fixed. Geometric ray errors are distinct from
the diffraction point-spread function and from measured optical quality.
